\section{Evaluation}\label{sec:experiments}
\input{tables/table1_comparison.tex}
\begin{figure}[t]
\centering
\includegraphics[width=\columnwidth]{figures/anchor_controlled_bounds.pdf}
\caption{Controlled four-leaf trees (inset: sibling pairs and references). (a) Increasing sibling-range overlap $d$ forces displacement; the width bound of Theorem~\ref{thm:width} (circles) equals the optimum over all 24 leaf orders. (b) At $d=24\%$, the 14\% displacement disappears exactly when the merge budget reaches half the parent gap $h$, for all four gaps, as Theorem~\ref{thm:chain} predicts. Percentages refer to a unit canvas or a unit scalar range.}
\label{fig:controlled}
\end{figure}
\begin{figure*}[t]
\centering
\includegraphics[width=\textwidth]{figures/plane_compact.pdf}
\caption{Position--topology plane of the bound suite, five planar and two volumetric sequences. Each map is one point: its worst merge error over frames (horizontal, square-root scale) and its mean order error (vertical). No map can enter the grey region below the bound (Theorem~\ref{thm:chain} and the flattened bound; Theorem~\ref{thm:order} at zero merge error). TMTM, ST-MTM and OLO keep the tree and sit at zero merge error far above the bound; on Storms 76--79\% of their order error is not forced by the tree. Space-filling curves and the projection follow space and break the tree. AnchorMTM (blue, $\varepsilon$ from 0 to 100\%) runs along the edge; on Storms, OLR, PM$_{2.5}$ and Ring, enumerating all leaf orders shows that it is the exact optimum at every tested budget except one Storms frame. TMTM, ST-MTM and OLO are shown in the mirror that agrees better with the reference.}
\label{fig:plane}
\end{figure*}
We ask four questions. \textbf{Q1:} are the bounds tight? \textbf{Q2:} how much of the positional error of existing maps does the tree force? \textbf{Q3:} at the same topological guarantee, does AnchorMTM place features closer to their references, and does it keep its budget? \textbf{Q4:} what does a budget buy, frame by frame? Section~\ref{sec:reading} then reads the maps.

\subsection{Data and Protocol}\label{sec:data}
\input{tables/table_data.tex}
Table~\ref{tab:data} lists two groups of sequences. The \emph{bound suite} comes with the saved trees of an earlier benchmark, prepared with one extractor, matching rule and reference estimator and with one default configuration; on it we evaluate the bounds, all leaf orders of small frames and every compared map frame by frame (Q1, Q2). The \emph{layout suite} reconstructs the four planar inputs of the ST-MTM benchmark~\cite{STMTM:2026}, so that AnchorMTM is compared on the data of the closest prior work, and adds the two TMTM volumes~\cite{Koepp:2023:TMTM,Weinkauf:SquareCylinder} (Q3, Q4). Ring, Storms, Wildfires and the volumes occur in both suites with different preprocessing (Supplementary Material). Author-processed arrays are unavailable, so the layout suite is a disclosed reconstruction, with the ST-MTM persistence numbers read as per-frame range fractions (0.1, 0.1, 15 and 10.5\%), not an identical-input reproduction.

\paragraph{Measures.} NMLE is the attained merge error $\dtop$ (Eq.~\ref{eq:dtop}) of the rendered profile divided by the scalar range, and NROE the order error of Eq.~\ref{eq:nroe}. NPE is the worst displacement $\max_i|x_i^t-q_i^t|$ of a frame divided by the physical extent and averaged over frames. SNS is the scale-normalized stress between support-centroid and anchor distances~\cite{Smelser:2025:Scale,STMTM:2026}, and TW measures neighborhoods~\cite{Venna:2001:NP} with $k=(1,3,5,3)$, using only frames with $0<k<n_t/2$ (64/64, 1/40, 4/32 and 11/744). The SNS and TW of TMTM and ST-MTM in Table~\ref{tab:quantitative-final} are quoted from ST-MTM (Temporal mode, authors' inputs) and are not ranked against local values. All other baseline values are measured on our implementations of the published methods with the same trees, identities and references, each in the global mirror that agrees better with the reference; Hilbert and Morton curves keep a fixed grid orientation~\cite{Franke:2021:1DP}. The volumetric table uses the $+x$ axis chosen by an offline displacement rule and stored, tuned configurations (Supplementary Table~S4), whereas Figs.~\ref{fig:volume-gallery} and~\ref{fig:budget-distribution} use the default configuration of the bound suite.

\paragraph{Modes.} Exact uses $\varepsilon=0$. Balanced selects a positive budget from 1--10\% on the first chronological quarter by $J=0.5\,\mathrm{NPE}_{norm}+0.5\,\mathrm{SNS}_{norm}$ over candidates within their budget, prefers smaller budgets within 0.01 of the minimum, and is frozen before the full runs: 1, 1, 5 and 10\% (Supplementary Fig.~S4). Earlier outputs had been observed, so this is not a blind test.

\subsection{Q1: Are the Bounds Tight?}\label{sec:bounds-eval}
\emph{Controlled trees.} We fix two sibling pairs $\{A,B\}$ and $\{C,D\}$ that merge at $0.4$ below a root at $1$, with leaf values $0$ and references $(0.2,0.5,0.5-d,0.8-d)$ on a unit canvas, and vary the signed overlap $d$ from $-0.12$ to $0.30$. Across 43 overlaps and four widths, we compare the bounds with all 24 leaf orders, each placed by an independent minimax LP. Topology alone forces exactly $\nu^*=\max(0,d)/2$. With widths $0.06$, gap $0.01$ and $\rho=0.5$, the width bound, our search and the exhaustive optimum coincide at $\tau=\max(0,(d+0.04)/2)$ (Fig.~\ref{fig:controlled}a): 12\% for points and 14\% for intervals at $d=0.24$. A merge budget removes the conflict exactly where Theorem~\ref{thm:chain} admits interleaving, at $2\varepsilon=h$ for all four parent gaps $h$ (Fig.~\ref{fig:controlled}b).

\emph{Real trees.} In Exact Mode, the interval cost $\tau$ of the selected order equals $\mathcal{B}_w$ on 606 of the 614 planar frames of the bound suite, so on these frames no exact-topology order has a smaller worst displacement; at $\varepsilon=10\%$ the bounded-topology bound is met on 596. On Storms, OLR, PM$_{2.5}$ and Ring every frame has at most eight leaves, and enumerating all leaf orders shows that the AnchorMTM order reaches the least order error at every tested budget, except one Storms frame at $\varepsilon=31.6\%$ (1.9 points). Volumes are harder: the width bound is met on 81 of 102 Cylinder and 27 of 201 Tangaroa frames, and the remaining gap may lie in the bound or in the heuristic search.

\begin{figure*}[t]
\centering
\includegraphics[width=\textwidth]{figures/volume_gallery.pdf}
\caption{Where each method draws the features of the two volumes. Each dot is one feature in one frame at its position on the axis (fraction of the canvas), coloured by its signed displacement from its reference; grey dots lie on their references. The first column shows the references, i.e., what a map read as location should show. All methods use the same trees and features; TMTM and ST-MTM are implementations of the published methods, each in the mirror that agrees better with the reference. AnchorMTM uses the default configuration of the bound suite with the matched-motion direction (Tangaroa) or the streamwise $x$ axis (Cylinder). Below each panel: mean and mean worst displacement over frames and, for a positive budget, the attained merge error.}
\label{fig:volume-gallery}
\end{figure*}
\subsection{Q2: How Much of the Error Does the Tree Force?}
Theorem~\ref{thm:order} needs no layout, so it measures every exact-topology map, and midpoint ordering attains $\nu^*$ on the same tree. On the seven bound sequences, 53--95\% of the order error of TMTM and 31--95\% of that of ST-MTM lies above $\nu^*$ and is therefore chosen by their order rules (Supplementary Table~S2); OLO, which only optimizes adjacency within the tree, behaves alike (32--95\%). Figure~\ref{fig:plane} places every map in the position--topology plane. Maps that keep the tree sit at zero merge error far above the bound; the projection reaches zero order error only with merge errors of 5.6--59\%, and the space-filling curves also break the tree (8--40\%); AnchorMTM runs along the edge between them. The edge shows what a budget releases: where conflicting subtrees merge just below their parent, 1\% halves the forced order error on Wildfires (5.80\% to 2.77\%) and almost removes it on Cylinder (12.24\% to 0.23\%), whereas the deep conflicts of Tangaroa keep it above 7\% up to $\varepsilon=5\%$. The forced cost also depends on the direction and can be computed before any map is drawn: over 400 directions it spans 4.70--10.14\% on Tangaroa and 5.49--28.15\% on Cylinder (Supplementary Fig.~S6).

\begin{figure*}[t]
\centering
\includegraphics[width=\textwidth]{figures/budget_distribution_2d3d.pdf}
\caption{Per-frame distribution of AnchorMTM over the merge budget $\varepsilon$ on complete sequences: three planar inputs of the layout suite with the configuration of Table~\ref{tab:quantitative-final}, and the two volumes with the default configuration. Top: worst displacement per frame (square-root scale; boxes: quartiles; whiskers: 5th--95th percentiles; dots: other frames; diamonds: means, the NPE of the tables). Bottom: attained merge error per frame with its budget (grey dash); no frame exceeds it. On the planar inputs a budget shortens the upper tail; in the volumes almost every frame uses it and the whole distribution moves. Gaussian is identical at every budget and omitted; Ring's selected 1\% and larger budgets appear in Supplementary Fig.~S1.}
\label{fig:budget-distribution}
\end{figure*}
\subsection{Q3: Closer at the Same Guarantee?}
\input{tables/table_volume.tex}
\emph{Planar inputs.} At the same zero merge error, Exact lowers NPE relative to TMTM and ST-MTM by 73--96\% on all four inputs of Table~\ref{tab:quantitative-final}, and it is also below the Hilbert and Morton curves, which break the tree. Balanced, with a budget of at most 10\%, lowers NPE relative to Exact by a further 16.2\% on Ring, 11.6\% on Wildfires and 64.9\% on Storms; Gaussian is unchanged. This supports alignment to a chosen reference, not general quality. ST-MTM reports higher TW on all four inputs, and we do not rank its published SNS and TW, which come from the authors' inputs and coverage, against local values. Against the curves, Balanced improves SNS and TW on Gaussian, Wildfires and Storms, but Hilbert is better on Ring (SNS 0.1252 and TW 0.9048 versus 0.1334 and 0.8571, with TW from one eligible frame), and Balanced raises SNS on Wildfires from 0.0426 to 0.0447.

\emph{Reading motion.} Positions are one reading; another asks whether two features that approach or separate in space also do so on the map. On the four planar bound sequences with such events (Storms, Wildfires, OLR and Ring; 56--124 events in which the distance between two tracked features changes by more than 2\% of the domain diagonal within two frames), an anchored exact-topology layout misreads 5--32\% of the events, TMTM 39--65\% and ST-MTM 21--41\%, and neither baseline misreads fewer events on any sequence. The task involves relative distances, which ST-MTM fits directly, and relates to the motion term of our placement, so it supports reading motion rather than providing independent evidence.

\emph{Volumes.} With the tuned configurations and the $+x$ reference of Table~\ref{tab:volume-comparison}, Exact lowers NPE at zero merge error from 0.477 to 0.177 on Cylinder and from 0.387 to 0.129 on Tangaroa, and Balanced improves all three metrics over TMTM; Tangaroa's 33.33\% budget and geometry-dominated weights limit attribution. Figure~\ref{fig:volume-gallery} shows the positions behind such numbers under the default configuration. TMTM and ST-MTM gather the Cylinder vortices near the middle of the axis (mean worst displacement 47.7\% and 51.6\%), whereas Exact already reproduces the shedding pattern of the references (9.2\%) and a 5\% budget removes almost all of the rest (1.2\%). Tangaroa's conflicts are deep: Exact lowers 30.6\% (TMTM) and 34.2\% (ST-MTM) to 11.2\%, and the pattern of the references appears near $\varepsilon=20\%$ (1.9\%).

\emph{Budgets.} Exact has zero measured merge error on every input. Balanced attains at most 0, 0.982, 4.421 and 9.985\%, within its budgets, and all runs pass continuous and raster checks without collapsed anchors; Ring needs checked joint integer rounding in 26 Exact and 21 Balanced frames. In the volumetric sweeps, all 1,515 rendered frames were rechecked against the fixed-leaf merge matrices without a violation. The checks cover diagonal extrema, not every augmented internal label.

\begin{figure*}[t]
\centering
\includegraphics[width=\textwidth]{figures/anchor_storms_budget_maps.pdf}
\caption{Storms parameter maps and method comparison. AnchorMTM spans exact topology ($\varepsilon=0$) and bounded topology (2, 5, 10 and 20\%), with the original trajectories and common local crops. TMTM and ST-MTM maps use the same scalar fields and augmented trees through local implementations. Spatial awareness and topology budget indicate the respective conceptual improvements. All maps show frames 355--460 with the original blue--orange scale, $[-35,35]$ hPa. Module connections summarize design differences rather than exact algebraic equivalence.}
\label{fig:storms-budget-profiles}
\end{figure*}
\subsection{Q4: What Does a Budget Buy?}
Figure~\ref{fig:budget-distribution} shows the budget frame by frame. On the planar inputs it shortens the upper tail rather than shifting every frame. On Storms the median per-frame displacement moves only from 1.10\% to 0.84\% at $\varepsilon=10\%$, while the 90th percentile falls from 10.7\% to 2.5\% and the worst frame from 44.8\% to 11.5\%; 482 of the 744 frames keep zero merge error, and those that use the allowance form a few runs of consecutive frames (Supplementary Fig.~S2). On Wildfires the same 7 of 32 frames use the allowance at every budget, and on Ring the median falls from 3.9\% to 1.0\% while the last frame stays at 10.0\% until the budget reaches 63\%. In the volumes almost every frame contains a releasable conflict (200 of 201 Tangaroa and 98 of 102 Cylinder frames use the allowance), so the whole distribution moves: the Cylinder median falls from 6.8\% to 1.0\% already at 2\%, and the Tangaroa median from 11.3\% to 1.2\% at 20\%. A budget limits merge error, not displacement, so a positive budget need not change the map, as on Gaussian.

\subsection{Reading the Maps}\label{sec:reading}
The upper row of Fig.~\ref{fig:storms-budget-profiles} shows a qualitative Storms archive, prepared separately from Table~\ref{tab:quantitative-final}, at budgets from 0 to 20\%. The 0, 2 and 5\% layouts overlap locally, while 10 and 20\% move the selected storm toward its fixed reference, so its track can be read as motion along the reference direction. The lower row shows TMTM and ST-MTM maps recomputed on the same 744 fields and augmented trees, in the same window and colours: TMTM aligns subtrees by temporal overlap and ST-MTM places them by relative distance, so a position on their axes does not correspond to a location of the storm. Both baselines are local implementations, not author binaries; genuine multiway saddles are handled by equal-level separators (TMTM) or an equal-level binary refinement whose leaf merge matrices are checked (ST-MTM).

\paragraph{Parameters and runtime.}
One cold Balanced run takes 0.5--63\,s on an Apple M4, of which ordering and placement take 0.15--3.3\,s; the rest is preprocessing, tree extraction, matching, rasterization and checks. Supplementary Tables~S3 and~S4 list all settings, and Supplementary Fig.~S8 shows a joint guidance--topology sweep on Ring with source snapshots.

